\chapter{Фредгольмовость составного дефекта и индекс Каллиаса}

Предыдущая проверка построила конечный классический профиль $(Q,A_Q)$. Теперь
проверяется следующий необходимый признак материи: существует ли на этом
фоне топологически защищённая локализованная фермионная мода.

Для оператора типа Каллиаса необходимы два раздельных условия:
\begin{enumerate}
 \item матричная масса должна быть обратима вне компактной области;
 \item положительное собственное расслоение массы на сфере бесконечности
 должно иметь ненулевой топологический класс.
\end{enumerate}
Первое условие обеспечивает фредгольмовость, второе --- ненулевой индекс.
См. DOI \path{10.1007/BF01202525}, DOI
\path{10.1007/BF01461005} и \path{arXiv:2108.06368}.

\section{Естественная масса на векторном носителе}

Самый экономный кандидат использует само поле порядка:
\begin{equation}
 \Phi_Q=\frac Q\Delta.
\end{equation}
На бесконечности его спектр равен
\begin{equation}
 \Spec\Phi_Q=\left\{\frac23,-\frac13,-\frac13\right\},
\end{equation}
поэтому масса равномерно обратима. В аналитическом смысле оператор
$D_{A_Q}+\Phi_Q$ является допустимым фредгольмовым кандидатом.

Но положительная собственная линия порождена вещественным директором
$n$. После выбора ориентированного подъёма она имеет глобальный
ненулевой вещественный раздел $n$; её комплексификация топологически
тривиальна:
\begin{equation}
 c_1\bigl(\mathbb C n\bigr)=0.
 \label{eq:v6-vector-callias-zero-chern}
\end{equation}
Прямой численный интеграл проекторной кривизны дал ноль с точностью
$10^{-12}$.

\begin{theorem}[Фредгольмовость без защищённой моды]
Естественная трёхкомпонентная масса $\Phi_Q$ обратима на бесконечности, но
её положительное собственное расслоение имеет нулевой первый класс Черна.
Поэтому индекс Каллиаса равен нулю.
\end{theorem}

Составная связность исправила энергию, но не изменила тип внутреннего
носителя. Нулевая мода на векторной ветви может возникнуть случайно, но не
защищена топологическим индексом.

\section{Масса спинорного накрытия}

Для ориентированного подъёма существует другой стандартный оператор:
\begin{equation}
 \Phi_{\rm sp}=n\cdot\sigma
\end{equation}
на комплексном носителе ранга два. Его положительный проектор
\begin{equation}
 P_+(n)=\frac{I_2+n\cdot\sigma}{2}
\end{equation}
задаёт хопфову линию $L$. Численный интеграл дал
\begin{equation}
 c_1(L)=1.000016\ldots,
\end{equation}
то есть единицу в пределах дискретизации. Минимальная ориентированная
ветвь имеет индекс по модулю один.

После тензорного умножения на коэффициентный проектор $q_0$ ранга
пятнадцать получаем
\begin{equation}
 c_1(L\otimes q_0)=15,
 \qquad
 |\operatorname{ind}\mathcal D_{\rm Callias}|=15.
 \label{eq:v6-spin-cover-callias-fifteen}
\end{equation}
Real-сопряжённая ветвь имеет противоположную ориентацию и индекс минус
пятнадцать.

Это совпадает с топологической классификацией дефекта $+15/-15$, но
требует комплексного носителя ранга два, на котором действуют матрицы
Паули независимо от пространственного клиффордова действия.

\section{Почему конечный родитель пока не даёт эту массу}

Том V исчерпывающе проверил внутренние кандидаты. Прямое нечётное поле
имеет тип
\begin{equation}
 T:\mathbb C^{15}\longrightarrow\mathbb C^{20}.
\end{equation}
Даже при максимальном ранге оно оставляет коядро размерности пять. Его
самосопряжённое удвоение имеет пять нулей на бесконечности и не
удовлетворяет условию Каллиаса.

Составная связность $A_Q$ действует на пространственных производных и не
меняет эти размерности. Она не превращает прямоугольный $20\times15$
угол в обратимый оператор спинорного накрытия ранга два.

Условные квадратные носители $15\times15$, $300\times300$ и
$150\times150$ требуют соответственно ручной редукции, скрытого удвоения
или нового комплексно-линейного спаривания. Поэтому конечный
$H_{15}/M_{35}$ всё ещё не выводит массу $n\cdot\sigma$.

\section{Совпадение с индексом Тёплица}

Здесь возникает важная положительная зацепка. Том V независимо построил
вещественный унитарий степени семь и его граница Тёплица:
\begin{equation}
 \operatorname{ind}T_{V_+}=-15,
 \qquad
 \operatorname{ind}T_{V_-}=+15.
 \label{eq:v6-toeplitz-callias-integer-match}
\end{equation}
Таким образом, три вычисления дают одно и то же абсолютное число:
\begin{equation}
 \boxed{
 15_{\rm spin\ cover}
 =15_{\rm Callias\ candidate}
 =15_{\rm Toeplitz\ boundary}.}
\end{equation}

Но равенство чисел ещё не является сравнительной теоремой. цикл Тёплица
живёт на $\ell^2(\mathbb Z)$ и пространстве Харди; новый оператор должен
жить на пространственной трёхмерной геометрии дефекта. Из стабильного
индекса Тёплица нельзя автоматически заключить существование
локализованных пространственных нулевых мод.

\begin{proposition}[Стабильный индекс найден, локализация не доказана]
Проект уже содержит полный Real-класс $\pm15$ в явном фредгольмовом
представителе Тёплица. Кандидат спинорного накрытия составного ежа имеет тот же
граничный индекс. Остаётся построить каспаровское сравнение Каллиаса--Тёплица
сравнительное отображение, сохраняющее пространственную локализацию.
\end{proposition}

\section{Следующий различающий тест}

Следующая проверка должна сравнить два граничного механизма:
\begin{enumerate}
 \item хопфов проектор $P_+(n)$ на $S^2_\infty$ и индекс Каллиаса;
 \item Real-унитарий $V_{15}(z)$, расширение Тёплица и индекс Харди.
\end{enumerate}
Нужно либо предъявить общий KK-произведённый класс и локализующий
оператор, либо доказать, что совпадение $15$ является только стабильным
числовым соответствием.

\begin{nogocriterion}[Индекс без носителя не локализует фермион]
Совпадение целых индексов Каллиаса и Тёплица не позволяет объявить
фермионную моду физической, пока не построен один оператор или явная
эквивалентность его пространственного представителя и представителя Харди.
\end{nogocriterion}

\begin{protocol}
\begin{enumerate}
 \item естественная масса $Q/\Delta$: \textbf{обратима на бесконечности};
 \item её индекс: \textbf{ноль};
 \item масса спинорного накрытия $n\cdot\sigma$: \textbf{индекс один};
 \item коэффициентный класс: \textbf{$+15/-15$};
 \item комплексный носитель ранга два из конечного родителя: \textbf{нет};
 \item прямой угол $20\times15$: \textbf{не имеет щели};
 \item вещественный индекс Тёплица: \textbf{$-15/+15$};
 \item совпадение абсолютных индексов: \textbf{да};
 \item пространственная локализация класса Тёплица: \textbf{не доказана};
 \item следующая проверка: \textbf{сравнение Каллиаса--Тёплица};
 \item локализованная фермионная материя: \textbf{не закрыта}.
\end{enumerate}
\end{protocol}

Машинный сертификат сохранён в
\path{s2t_v6_composite_connection_callias_fredholm_gate_results.json}.